\section{Conclusion}
\label{sec:conclusion}

We presented \sys, a cross-stage pruning framework for disaggregated prefill--decode agentic serving.
On Qwen3 models and reasoning workloads, prefill-side admission reduces computed prefill by $23\%$ at matched committed-answer accuracy, a $128$-token probe restores MATH-500 coverage to $20/32$ at $0.69\times$ the full-fan-out wall-clock, and spine retention reduces median peak KV usage to $0.72\times$ that of APC.
These results show the benefit of filtering branches at the earliest informed stage and of retaining only the committed spine.

\paragraph{Limitations.}
\label{sec:limitations}
\sys currently relies on lightweight heuristic scores for admission and probing, evaluated primarily on reasoning-oriented branching workloads.
Although the framework is agnostic to the scoring function, the separation observed in our experiments may not generalize directly to heterogeneous agent traces, tool-use trajectories, or less structured residuals.
Evaluating learned or workload-adaptive scoring policies on large-scale production traces remains an important direction for future work.
